\documentclass[a4paper,12pt]{article}
\usepackage{xeCJK}          % 中文支持
\usepackage{fontspec}       % 英文/数学字体
\usepackage{amsmath, amssymb} % 数学公式
\usepackage{graphicx}       % 插入图片
\usepackage{hyperref}       % 目录超链接
\usepackage{geometry}       % 页面布局
\usepackage{bm}             % 粗体
\usepackage{xcolor}         % 颜色
\usepackage{tabularx}       % 表格环境
\usepackage{tikz}           % TikZ 绘制主对角线斜线
\usepackage{tcolorbox}
\usepackage{xstring}
\usepackage{pgfplots}
\usepackage{enumitem}
\usepackage{pifont}
\usepackage{ulem}
\usepackage{mathtools}
\usepackage{dsfont}
\usepackage{extarrows}
\pgfplotsset{compat=1.18}
\geometry{left=3cm,right=3cm,top=3cm,bottom=3cm}

\definecolor{customgreen}{rgb}{0.2,0.6,0.3}
\newcommand{\gmath}[1]{\textcolor{customgreen}{#1}}

% 经典
\definecolor{classicgreen}{rgb}{0.1,0.5,0.3}
\newtcbox{\classic}{
    on line,
    colback=classicgreen!10,
    colframe=classicgreen,
    boxrule=0.8pt,
    arc=3pt,
    boxsep=1pt,
    left=4pt,
    right=4pt,
    top=2pt,
    bottom=2pt
}

% 抽离颜色和尺寸参数
\newcommand{\analysisTitleColor}{green!50!black}
\newcommand{\analysisBackColor}{white}
\newcommand{\analysisBoxRule}{0.8pt}
\newcommand{\analysisArc}{3pt}
\newcommand{\analysisPadding}{6pt}

% 定义 tcolorbox
\newtcolorbox{analysisbox}[1][]{
    title=\IfStrEq{#1}{}{\textbf{解析}}{#1}, % 如果传参为空则使用“解析”
    colback=\analysisBackColor,
    colframe=\analysisTitleColor,
    boxrule=\analysisBoxRule,
    arc=\analysisArc,
    left=\analysisPadding,
    right=\analysisPadding,
    top=4pt,
    bottom=4pt
}

\newlist{circlenum}{enumerate}{1}
\setlist[circlenum]{
    label=\ding{\numexpr171+\arabic*},
    leftmargin=2.2em,
    itemsep=0.6em,      % item 之间的距离（主要）
    topsep=0.6em,       % 列表与上下正文的距离
    parsep=0.3em,       % item 内段落的间距
    partopsep=0.3em     % 列表前后额外间距
}

\newcommand{\blueuline}[1]{{\color{blue}\uline{\color{black}{#1}}}}

% =========================
% 字体设置
% =========================
\setmainfont{Times New Roman}
\setsansfont{Helvetica Neue}
\setmonofont{Menlo}
\setCJKmainfont{PingFang SC}

% =========================
% 图形路径（可调整）
% =========================
\graphicspath{{./assets/}}

% =========================
% 文档开始
% =========================
\begin{document}

%    \title{Template}
%    \author{Bowen}
%    \date{\today}
%    \maketitle
%    \tableofcontents
%    \newpage


% =========================

    \section{一元函数积分学的应用 - 积分等式与积分不等式}

    \subsection{积分等式}

    \begin{enumerate}
        \item 用中值定理
        \item 用夹逼准则
        \item 用积分法
    \end{enumerate}

    \subsection{积分不等式}

    \begin{enumerate}
        \item 用函数的单调性：首先将某一积分限（通常取上限）变量化，然后移项构造辅助函数，由辅助函数的单调性来证明不等式，此方法多用于所给条件为\textbf{\blueuline{$f(x)$在$[a, b]$上连续}}的情形
        \item 用拉格朗日中值定理：此方法多用于所给条件为\textbf{\blueuline{$f(x)$一阶可导}}且某一端点值比较简单（甚至为$0$）的题目
        \item 用泰勒公式：此方法多用于所给条件为\textbf{\blueuline{$f(x)$二阶可导}}且题中有简单函数值（甚至为$0$）的题目
        \item 用积分法
        \item 用牛顿-莱布尼茨公式
    \end{enumerate}

    \subsection{基础概念}

    \begin{enumerate}
    \end{enumerate}

    \subsection{结论}

    \begin{enumerate}
    \end{enumerate}

    \subsection{定理}

    \begin{enumerate}
    \end{enumerate}

    \subsection{运算}

    \begin{enumerate}

    \end{enumerate}

    \subsection{公式}

    \begin{enumerate}

    \end{enumerate}

    \subsection{方法总结}

    \begin{enumerate}

    \end{enumerate}

    \subsection{条件转换思路}

    \begin{enumerate}

    \end{enumerate}

    \subsection{理解}

    \begin{enumerate}
    \end{enumerate}

\end{document}
